El proceso actual de gestión y seguimiento de tesis y proyectos de grado en la FICH carece de mecanismos sistematizados que garanticen eficiencia en los flujos de trabajo, control sobre el estado de cada proyecto, calidad en la producción intelectual y un respaldo adecuado del conocimiento generado. Actualmente no existe una plataforma unificada que integre el registro de manuscritos, la asignación de tutores y tribunales, el flujo de revisiones y retroalimentación académica, la consulta pública del estado de cada proyecto, la gestión de pagos, ni la generación de reportes y estadísticas. Esta ausencia deriva en retrasos, pérdida de información, falta de trazabilidad y dificultades para que estudiantes, docentes y autoridades den seguimiento oportuno a los procesos de titulación.
